% ==============================================================================
% Appendix C — Description of datasets
% ==============================================================================

\chapter{Description of datasets}
\label{app:datasets}

This section documents the main external datasets used across the spatial
proxy construction pipeline.

\subsection{European boundaries}
\label{app:boundaries}

The European boundaries are used to clip the datasets to the European domain. The boundaries are provided by the European Comission and are available at the NUTS-2 level.
They are provided by Geographic Information System of the (European) Commission.
The dataset is used across the different proxy sectors for clipping, aggregation and reporting
We present in ~\ref{fig:GRC-NUTS2} the different regions of greece from the dataset.

\begin{figure}[htbp]
    \centering
    \includegraphics[width=0.8\textwidth]{figures/Dataset/greece_nuts2.png}
    \caption{NUTS-2 regions of Greece as used for spatial domain definition and reporting.}
    \label{fig:GRC-NUTS2}
\end{figure}

The file can be accessed at this adress : 
\hyperref[app:boundaries]{https://gisco-services.ec.europa.eu/distribution/v2/nuts/gpkg/NUTS_RG_03M_2021_3035.gpkg}
On the \hyperref[GISCO]{https://gisco-services.ec.europa.eu/distribution/v2/nuts/nuts-2021-files.html} website. No account is required.

\subsection{CORINE Land Cover}

The CORINE Land Cover is used to identify the different land cover types across the European domain.
It is provided by the European Environment Agency (EEA) and is available at the 100 m resolution.
It is directly used by CAMS for different sectors to identify land type and thus emissions from different sectors. Especially, the industrial land cover and and agriculture. \ref{}
And it thus higly relevant to downscale emissions from different sectors to the fine-scale pixels.
Below we present a table ~\ref{tab:CLC-codes} of CORINE land cover types and their corresponding codes.

\begin{table}[htbp]
    \centering
    \begin{tabular}{|c|l|c|l|}
        \hline
        \textbf{Code} & \textbf{Description} & \textbf{Code} & \textbf{Description} \\
        \hline
        111 & Continuous urban fabric & 321 & Natural grasslands \\
        112 & Discontinuous urban fabric & 322 & Moors and heathland \\
        121 & Industrial or commercial units & 323 & Sclerophyllous vegetation \\
        122 & Road and rail networks and associated land & 324 & Transitional woodland-shrub \\
        123 & Port areas & 331 & Beaches dunes sands \\
        124 & Airports & 332 & Bare rocks \\
        131 & Mineral extraction sites & 333 & Sparsely vegetated areas \\
        132 & Dump sites & 334 & Burnt areas \\
        133 & Construction sites & 335 & Glaciers and perpetual snow \\
        141 & Green urban areas & 411 & Inland marshes \\
        142 & Sport and leisure facilities & 412 & Peat bogs \\
        211 & Non-irrigated arable land & 421 & Salt marshes \\
        212 & Permanently irrigated land & 422 & Salines \\
        213 & Rice fields & 423 & Intertidal flats \\
        221 & Vineyards & 511 & Water courses \\
        222 & Fruit trees and berry plantations & 512 & Water bodies \\
        223 & Olive groves & 521 & Coastal lagoons \\
        231 & Pastures & 522 & Estuaries \\
        241 & Annual crops associated with permanent crops & 523 & Sea and ocean \\
        242 & Complex cultivation patterns & 999 & NODATA \\
        243 & Land principally occupied by agriculture & & \\
        & with significant areas of natural vegetation & & \\
        244 & Agro-forestry areas & & \\
   
        311 & Broad-leaved forest & & \\
        312 & Coniferous forest & & \\
        313 & Mixed forest & & \\
        \hline
    \end{tabular}
    \caption{CORINE land cover types and their corresponding codes (\texttt{CLC2018\_CLC2018\_V2018\_20\_QGIS.txt}).}
    \label{tab:CLC-codes}
\end{table}

The dataset is used across the different proxy sectors for clipping, aggregation and reporting



\subsubsection*{European open power-plant register}
\label{app:powerplant_register}

The point-source Public Power proxy uses a European open register of
power-generation units. The dataset is provided at unit level rather than only
at plant level, and includes geographic coordinates, country, operating
status, installed capacity, and generation type. It combines information from
ENTSO-E and other open sources into a harmonised European table, but should be
understood as an open reference register rather than a complete census of all
facilities.

Only combustion-based public-power units are retained for the proxy, because
the objective is to localise stack-related emissions from electricity and heat
production. Hydro, wind, solar, and other non-combustion generation types are
excluded. The main fields used are therefore the unit coordinates, the country
identifier, and the generation-type classification. The dataset is used only
as a localisation support for point emissions; it does not provide plant-level
emissions directly.

\subsubsection*{LUCAS (Land Use/Cover Area frame Survey)}
\label{app:lucas}
LUCAS is a pan-European in-situ field survey conducted by Eurostat at a stratified systematic grid of approximately 1~km spacing, with approximately 270\,000--340\,000 field observations per survey wave \citep{lucas_dataset}. The 2022 wave, used throughout this work, covers EU-27 member states plus selected candidate countries. Each survey point records:
\begin{itemize}
  \item \textbf{Land cover} (\texttt{SURVEY\_LC1}): a hierarchical code identifying the dominant cover type (e.g.\ B11 = wheat, E10 = permanent grass, B17 = rice). Level-1 codes B* denote cropland, E* denote grassland and heath.
  \item \textbf{Land use} (\texttt{SURVEY\_LU1}): functional category (U111 = agriculture including animal production, U112 = fallow, U113 = kitchen garden, U120 = forestry).
  \item \textbf{Grazing} (\texttt{SURVEY\_GRAZING}): integer indicator (0 = no signs, 1 = grazing observed, 2 = no signs at time of observation).
  \item \textbf{INSPIRE patch fractions}: \texttt{SURVEY\_INSPIRE\_ARABLE} (arable fraction, \%), \texttt{SURVEY\_INSPIRE\_ORGCON} (organic deposit fraction, \%), \texttt{SURVEY\_INSPIRE\_PLCC4} (herbaceous cover, \%), \texttt{SURVEY\_INSPIRE\_PLCC7} (bare/unconsolidated, \%).
  \item \textbf{Crop residue management}: \texttt{SURVEY\_LM\_CROP\_RESID} (presence code) and \texttt{SURVEY\_LM\_CROP\_RESID\_PERC} (estimated residue \%).
\end{itemize}
LUCAS points are intersected with the CORINE Land Cover raster at each coordinate; only points falling on agricultural CORINE classes (raster codes 12--22) are retained in the proxy pipeline. The survey is the only systematic in-situ observational basis for inferring land-management intensity at sub-NUTS-2 resolution across all EU member states simultaneously.

\subsection{LUCAS 2018 Topsoil}
\label{app:lucas_topsoil}
\subsubsection*{Eurostat Agricultural Census Livestock Counts (C21)}
\label{app:c21}
The C21 dataset is Eurostat's Farm Structure Survey (FSS/SAPM) livestock count data, harmonised at NUTS-2 resolution and provided as a GeoPackage (\texttt{C21.gpkg}). It records the number of livestock heads (thousands) by species and NUTS-2 region, for reference years aligned to the EU Farm Structure Survey cycles (most recent: 2020). Species layers used in this work are bovines (total and dairy), pigs (under and over 50~kg), sheep, goats, and poultry. Counts are multiplied by IPCC-2006 species-specific emission factors to derive stage-2 census intensity weights $\omega_{n,c}$ for enteric fermentation (Section~\ref{sec:mu-rho-enteric}) and livestock housing (Section~\ref{sec:proxy-housing}). Column names within the GeoPackage are mapped via \texttt{c21\_census\_field\_map.json} to avoid hardcoding in Python.

\subsubsection*{GFED4.1s (Global Fire Emissions Database version 4.1 with small fires)}
\label{app:CGEF4.1s}
GFED4.1s \citep{vanderwerf2017,giglio2013} is a global monthly 0.25\textdegree{} gridded fire emissions product covering 1997--2015. It is based on MODIS burned area (Collection~5 MCD45/MCD64), active fire counts, and a biogeochemical model (CASA). For each month and pixel, the dataset provides:
\begin{itemize}
  \item \textbf{Total dry matter burned} (\texttt{emissions/MM/DM}), kg~DM~m$^{-2}$~month$^{-1}$.
  \item \textbf{Source-type fractions} (\texttt{emissions/MM/partitioning/DM\_AGRI} and five analogous fields), unitless $\in [0,1]$, summing to exactly 1.0 per pixel per month. The six types are: agricultural waste burning (\textsc{agri}), boreal forest (\textsc{borf}), deforestation (\textsc{defo}), peat (\textsc{peat}), savanna and grassland (\textsc{sava}), and temperate forest (\textsc{temf}).
  \item \textbf{Grid-cell area} (\texttt{ancill/grid\_cell\_area}), m$^2$.
  \item \textbf{Burned fraction} (\texttt{burned\_area/MM/burned\_fraction}), fraction of grid cell burned per month.
\end{itemize}
The annual files are in HDF5 format, located at \texttt{data/Agriculture/fire\_emissions\_v4\_R1\_1293/data/GFED4.1s\_<year>.hdf5}. The one-time preprocessing script \texttt{Agriculture/preprocess/build\_gfed41s\_agri\_mean.py} reads all 19 annual files, computes the multi-year mean agricultural DM per pixel (Eq.~\ref{eq:gfed-score}), and builds the pixel-to-NUTS-2 lookup table. Output artefacts are saved to \texttt{data/Agriculture/} as NumPy binary and Parquet files. The 19-year climatological mean (1997--2015) stabilises the spatial pattern against inter-annual fire anomalies: over the EU domain, the coefficient of variation of annual AGRI DM across years is approximately 34\,\%.

\subsubsection*{CEIP/NRD waste emissions}
\label{app:waste_ceip}

The waste downscaling workflow uses a country-level CEIP/NRD workbook of
waste-sector emissions with pollutant and subsector detail. The relevant
subsectors are 5A, 5B1, 5B2, 5C2, 5D1, 5D2, 5D3, and 5E, while the
incineration-like 5C1$\ast$ rows are retained only for bookkeeping and are
excluded from the default area-weighting logic. The main fields required by the
processing code are country identifier, year, pollutant, subsector code, and
reported total. These records are aggregated to derive country- and
pollutant-specific family shares for solid waste treatment, wastewater, and
residual diffuse waste. The workbook is therefore used as an apportionment
source for the proxy mixture, not as a gridded emissions product.

\subsubsection*{UWWTD agglomerations and treatment plants}
\label{app:waste_uwwtd}

The Urban Waste Water Treatment Directive (UWWTD) layers provide the main
infrastructure support for the wastewater family. Two complementary datasets are
used: agglomerations, which delineate the serviced wastewater population
footprint, and treatment plants, which locate the corresponding treatment
infrastructure. In the proxy workflow, agglomeration polygons provide diffuse
service-area support, while treatment plants act as local infrastructure
anchors. The datasets are used only as spatial indicators of wastewater
handling; they do not provide pollutant-resolved emissions per facility or per
agglomeration for direct downscaling.

\subsubsection*{GHSL population and settlement model}
\label{app:waste_ghsl}

The waste proxy workflow uses the GHSL population grid and the GHSL Degree of
Urbanisation / SMOD settlement classification as harmonised European settlement
support. The population layer provides a continuous indicator of human presence,
which is relevant both for domestic wastewater generation and for diffuse
waste-related activity. The SMOD layer distinguishes dense urban, peri-urban,
and rural settlement structures and is used mainly in the residual diffuse
waste family to emphasise rural and low-density settlement environments where
open waste burning and other diffuse waste processes are more plausible than in
large specialised facilities.

\subsubsection*{Copernicus imperviousness}
\label{app:waste_imperviousness}

The Copernicus High Resolution Layer on imperviousness is used as a built-surface
indicator complementing population and settlement classes. In the wastewater and
residual diffuse waste families, imperviousness acts as a fine-scale proxy for
sealed or strongly urbanised surfaces where waste-related activity is more
likely than in natural land-cover classes. The layer is used only as a relative
spatial support. Missing coverage is treated explicitly in the implementation
and recorded through diagnostic masks so that the weighting workflow remains
transparent where only partial tile coverage is available.

\subsubsection*{OpenStreetMap waste-related features}
\label{app:waste_osm}

OpenStreetMap data are used only as optional geometric refinement for the solid
waste family. Relevant tags include landfill- or waste-related land use and
amenity classes where such features can be extracted consistently from the
available OSM source. Because completeness and tagging practice vary across
space, OSM does not define the waste proxy on its own; instead, it refines the
CORINE-based support where compatible features are present. The workflow is
therefore robust to missing or partial OSM coverage, but the geometric detail is
better where the underlying OSM extract is more complete.
